\section{Risk Management}\label{sec:risk-management}
The working assessment contains twelve risks and proposed responses. It follows
the risk identification/analysis stages and continuous review cycle in the supplied course material~\cite{course-risk}. It is a desk-based consolidation
of the team's contribution; no facilitated workshop, complete nine-person vote
or completed mitigation is claimed. Ratings and named response ownership still
require team validation under \decisionref{Q-PM-01}.

\subsection{Approach and Assessment Scales}
Identify a present condition and its possible adverse consequence, keeping the
statement short and recording context/triggers separately. Review three
project-specific source groups: technical design/integration, delivery/capacity,
and stakeholder validation. This grouping organises the register; it is not a
claim that a full SEI taxonomy questionnaire was completed. An actual defect or
missed milestone enters the issue log, while its possible wider effects remain
risks. Communication accompanies identification, analysis, planning, tracking and
control throughout delivery.

The supplied probability bands are judgement categories, not measured frequencies:
Low below 40\%, Medium 40\% through 70\%, High above 70\%. Impact is measured
against the candidate threshold of success: Marginal permits success with limited
rework; Critical permits it only with substantial effort or adjustment;
Catastrophic prevents it within the agreed baseline. Timeframe means the time
until action is needed: Short within one month, Medium one to three months,
Long beyond three months. It is not the expected duration of an incident.
Confirm timeframes against the actual remaining project horizon.

Exposure uses the categorical lookup in Table~\ref{tab:risk-exposure}, not an
unexplained multiplication of ordinal labels. The matrix contains the current
provisional risk placements and is consistent with the register.
\Needspace{12\baselineskip}
{\compacttable
\begin{longtable}{L{0.16}L{0.28}L{0.28}L{0.28}}
\caption{Qualitative exposure matrix and provisional risk placement.}\label{tab:risk-exposure}\\
\toprule\rowcolor{TableHeader}
\textbf{Probability} & \textbf{Marginal impact} & \textbf{Critical impact} & \textbf{Catastrophic impact}\\\midrule\endfirsthead
\toprule\rowcolor{TableHeader}\textbf{Probability} & \textbf{Marginal impact} & \textbf{Critical impact} & \textbf{Catastrophic impact}\\\midrule\endhead
\bottomrule\endlastfoot
High & Medium: \identifier{RS-04}, \identifier{RS-12} & High: \identifier{RS-01}, \identifier{RS-07}, \identifier{RS-08}, \identifier{RS-09} & Critical: \identifier{RS-06}\\\midrule
Medium & Low: \identifier{RS-03}, \identifier{RS-05}, \identifier{RS-11} & Medium: \identifier{RS-02}, \identifier{RS-10} & High: none identified\\\midrule
Low & Low: none identified & Low: none identified & Medium: none identified\\
\end{longtable}}

\subsection{Candidate Threshold of Success}\label{sec:threshold-success}
The risk-assessment process uses a threshold of success (ToS):
a small set of measurable, time-bound conditions that must all be met. Table~\ref{tab:success-threshold} makes the
proposed conditions and evidence explicit. They do not yet constitute an agreed
ToS: acceptance principles and review roles are agreed, while final dates, named
owners, execution evidence and resource commitments remain to confirm. The intermediate submission milestone of 6 November 2026 and
the prototype presentation milestone of 10 November 2026 are distinct past
checkpoints; their dates do not by themselves demonstrate completion. The final delivery and
defence dates remain to be agreed under \decisionref{Q-PM-01}.

{\compacttable
\begin{longtable}{L{0.14}L{0.57}L{0.29}}
\caption{Candidate success conditions and remaining baseline inputs.}\label{tab:success-threshold}\\
\toprule\rowcolor{TableHeader}\textbf{ID / area} & \textbf{Condition and evidence at the agreed delivery} & \textbf{Input still to agree}\\\midrule\endfirsthead
\toprule\rowcolor{TableHeader}\textbf{ID / area} & \textbf{Condition and evidence at the agreed delivery} & \textbf{Input still to agree}\\\midrule\endhead
\multicolumn{3}{r}{\footnotesize\itshape Continued on next page}\\\endfoot
\bottomrule\endlastfoot
ToS-1\newline Scope & All criteria included in the final-delivery baseline have recorded acceptance evidence, including its field/document paths. & Execute criteria; name reviewers and confirm final date; \decisionref{Q-SCOPE-01}.\\\midrule
ToS-2\newline Schedule & Required software/prototype, repository, report and presentation are delivered by their agreed milestone dates, with dated submission evidence. & Final calendar and responsibilities; \decisionref{Q-PM-01}.\\\midrule
ToS-3\newline Quality & No known defect blocks a principal flow, permits unauthorised access, loses data, duplicates effects or produces incorrect totals. Included criteria have verification evidence; minor issues are documented and accepted in review. & Execute the agreed acceptance checks and record reviewers/results; Section~\ref{sec:validation}.\\\midrule
ToS-4\newline Field use & Selected phone/tablet/desktop tasks and prepared field flows pass the agreed offline/reconnection and usability criteria, with durable proposals, dependency validation and no silent conflict overwrite. & Record exact devices/browsers and execution evidence; Section~\ref{sec:nfr-validation}.\\\midrule
ToS-5\newline Resources & Actual effort and service expenditure remain within the approved effort/infrastructure budget or a recorded authorised revision. & Applicable per-member effort and service limits; Section~\ref{sec:project-planning}.\\
\end{longtable}}

\subsection{Risk Register}
Statements separate the present source of concern from its possible loss. Context,
trigger and proposed role owner accompany the short statement; probability, impact
and timeframe are inherited from the contribution pending team confirmation.
\emph{ToS} numbers identify potentially affected success conditions, not proof of
approved priorities. Re-rate impacts when the selected scope and evidence are known.
{\compacttable
\begin{longtable}{L{0.11}L{0.62}L{0.27}}
\caption{Working risk register and qualitative assessment.}\label{tab:risk-register}\\
\toprule\rowcolor{TableHeader}
\textbf{ID} & \textbf{Condition, consequence and trigger} & \textbf{Assessment / proposed owner}\\\midrule\endfirsthead
\toprule\rowcolor{TableHeader}
\textbf{ID} & \textbf{Condition, consequence and trigger} & \textbf{Assessment / proposed owner}\\\midrule\endhead
\multicolumn{3}{r}{\footnotesize\itshape Continued on next page}\\\endfoot
\bottomrule\endlastfoot
\phantomsection\label{risk:RS-01}\identifier{RS-01} & \textbf{Architecture validation.} The selected stack and candidate services still require an integration slice; incompatible choices may delay implementation or cause rework.\newline\textit{Source:} Technical / integration. \textit{ToS:} 1, 2, 5.\newline\textit{Trigger:} First integration slice blocked by an unresolved contract, dependency or service choice. & High / Critical / Short.\newline\textbf{Exposure: High.}\newline Backend/architecture lead.\\
\midrule
\phantomsection\label{risk:RS-02}\identifier{RS-02} & \textbf{Offline synchronisation.} Field workflows require durable queues, dependency ordering and concurrent updates; integration complexity may cause reconciliation failures or delay the agreed delivery.\newline\textit{Source:} Technical / integration. \textit{ToS:} 1, 3, 4.\newline\textit{Trigger:} Retries duplicate effects, pending work disappears, or stale revisions overwrite accepted data in tests. & Medium / Critical / Short.\newline\textbf{Exposure: Medium.}\newline Backend/synchronisation lead.\\
\midrule
\phantomsection\label{risk:RS-03}\identifier{RS-03} & \textbf{Outdoor interaction.} Phone/tablet field use is in scope, but representative outdoor validation is not evidenced; unsuitable controls may prevent efficient logging.\newline\textit{Source:} Technical / usability. \textit{ToS:} 1, 4.\newline\textit{Trigger:} Representative phone/tablet users cannot complete the selected tasks under agreed conditions. & Medium / Marginal / Medium.\newline\textbf{Exposure: Low.}\newline UI/usability lead.\\
\midrule
\phantomsection\label{risk:RS-04}\identifier{RS-04} & \textbf{Regression readiness.} The report contains planned checks but no verified CI/regression evidence; missing automated checks may allow defects to reach a milestone demonstration.\newline\textit{Source:} Technical / integration. \textit{ToS:} 2, 3.\newline\textit{Trigger:} A change reaches integration without passing the agreed build and core-rule checks. & High / Marginal / Short.\newline\textbf{Exposure: Medium.}\newline Testing/integration lead.\\
\midrule
\phantomsection\label{risk:RS-05}\identifier{RS-05} & \textbf{Contribution consistency.} Multiple contributors share client, server and rule contracts; inconsistent implementation or unreviewed generated code may create incompatible dependencies and rework.\newline\textit{Source:} Technical / integration. \textit{ToS:} 2, 3, 5.\newline\textit{Trigger:} An unreviewed dependency, incompatible interface or unexplained generated block enters a proposed merge. & Medium / Marginal / Short.\newline\textbf{Exposure: Low.}\newline Technical reviewer.\\
\midrule
\phantomsection\label{risk:RS-06}\identifier{RS-06} & \textbf{Academic capacity.} Nine students share project work with other academic commitments; reduced availability may leave insufficient effort for implementation and acceptance.\newline\textit{Source:} Delivery / capacity. \textit{ToS:} 1, 2, 5.\newline\textit{Trigger:} Forecast work exceeds available person-hours or an owner cannot meet a dependent task date. & High / Catastrophic / Short.\newline\textbf{Exposure: Critical.}\newline Project coordinator.\\
\midrule
\phantomsection\label{risk:RS-07}\identifier{RS-07} & \textbf{Scope expansion.} Some optional capabilities still need explicit scope decisions for the single final delivery; uncontrolled additions may displace essential development and validation work.\newline\textit{Source:} Delivery / scope. \textit{ToS:} 1, 2, 5.\newline\textit{Trigger:} An unestimated feature or excluded payroll/payment/mapping capability enters the delivery backlog. & High / Critical / Short.\newline\textbf{Exposure: High.}\newline Project coordinator.\\
\midrule
\phantomsection\label{risk:RS-08}\identifier{RS-08} & \textbf{Client clarification.} Ticket/document interpretation depends on client clarification through an intermediary; delayed answers may block correct implementation or acceptance.\newline\textit{Source:} Stakeholder validation. \textit{ToS:} 1, 2, 3.\newline\textit{Trigger:} A decision needed by a scheduled task has no validated answer or agreed fallback by its review date. & High / Critical / Short.\newline\textbf{Exposure: High.}\newline Client intermediary.\\
\midrule
\phantomsection\label{risk:RS-09}\identifier{RS-09} & \textbf{Task ownership.} Nine members contribute across connected deliverables; unclear ownership or stale tracking may create gaps, duplicated effort and uneven workload.\newline\textit{Source:} Delivery / coordination. \textit{ToS:} 1, 2, 5.\newline\textit{Trigger:} A deliverable lacks one accountable owner/date or completed work is absent from the shared tracking. & High / Critical / Short.\newline\textbf{Exposure: High.}\newline Project coordinator.\\
\midrule
\phantomsection\label{risk:RS-10}\identifier{RS-10} & \textbf{Integration friction.} Parallel client, API and schema changes share contracts; incompatible changes may break integration and consume the planned validation period.\newline\textit{Source:} Technical / integration. \textit{ToS:} 2, 3, 5.\newline\textit{Trigger:} A dependent branch cannot merge/build or a contract/migration change has no coordinated consumer update. & Medium / Critical / Short.\newline\textbf{Exposure: Medium.}\newline Integration lead.\\
\midrule
\phantomsection\label{risk:RS-11}\identifier{RS-11} & \textbf{Handover continuity.} Delivery tasks need continuing ownership despite changing availability; incomplete handovers may leave decisions or unfinished work without an accountable owner.\newline\textit{Source:} Delivery / coordination. \textit{ToS:} 1, 2.\newline\textit{Trigger:} A departing/unavailable owner has no documented handover, backup or accessible work. & Medium / Marginal / Medium.\newline\textbf{Exposure: Low.}\newline Project coordinator.\\
\midrule
\phantomsection\label{risk:RS-12}\identifier{RS-12} & \textbf{Representative validation.} Acceptance depends on representative users, devices and source examples that are not yet fully evidenced; synthetic scenarios alone may conceal unsuitable workflows.\newline\textit{Source:} Stakeholder validation. \textit{ToS:} 1, 3, 4.\newline\textit{Trigger:} Planned validation lacks a representative user, device or sanitised source example. & High / Marginal / Medium.\newline\textbf{Exposure: Medium.}\newline Client intermediary / usability lead.\\
\end{longtable}}
The assessment column lists probability / impact / timeframe, followed by exposure and proposed role owner.


\subsection{Priorities and Response Plans}
The source's six focus risks are retained as a provisional intervention list.
Exposure puts \identifier{RS-06} first and the four High-exposure risks next;
\identifier{RS-02} also receives early attention because offline behaviour is a
prerequisite for the field workflow. This explains its inclusion alongside, rather
than automatically above, other Medium-exposure risks. Reassess this tie-breaker
once the ToS and implementation plan are approved. A low vote never removes a
mandatory security/integrity check.

Multi-voting is optional. The supplied eight-column ballot does not establish a
complete nine-member vote. If repeated with four votes per participant, nine
participants would cast 36 votes across the twelve risk statements. Record actual
participation, discuss ties and discrepancies, and agree the ranking before treating
it as a team result.

{\compacttable
\begin{longtable}{L{0.14}L{0.23}L{0.63}}
\caption{Initial intervention priorities and affected project work.}\label{tab:risk-priority-links}\\
\toprule\rowcolor{TableHeader}\textbf{Risk} & \textbf{Response strategy} & \textbf{Why act / linked work}\\\midrule\endfirsthead
\toprule\rowcolor{TableHeader}\textbf{Risk} & \textbf{Response strategy} & \textbf{Why act / linked work}\\\midrule\endhead
\multicolumn{3}{r}{\footnotesize\itshape Continued on next page}\\\endfoot
\bottomrule\endlastfoot
\identifier{RS-06} & Reduction & Critical exposure affects deliverable capacity; reconcile actual/remaining person-hours and owner availability (Section~\ref{sec:project-planning}).\\\midrule
\identifier{RS-01} & Avoidance / reduction & High exposure affects implementation dependencies; validate \decisionref{Q-ARCH-01}/\decisionref{Q-ARCH-02} with the authenticated vertical slice.\\\midrule
\identifier{RS-07} & Avoidance & High exposure threatens scope/effort; apply \decisionref{Q-SCOPE-01} and the change process in Section~\ref{sec:requirements-change}.\\\midrule
\identifier{RS-08} & Reduction & High exposure blocks source/template acceptance; validate supplied examples and obtain missing clarification for \decisionref{Q-DA-02}, \decisionref{Q-DOC-01} and \decisionref{Q-DOC-02}.\\\midrule
\identifier{RS-09} & Reduction & High exposure threatens ownership and dependencies; reconcile task board, contribution review and hours evidence.\\\midrule
\identifier{RS-02} & Reduction & Medium exposure with core field dependency; implement/test \identifier{REL-01}--\identifier{REL-03} and \identifier{AOC-02}/\identifier{AOC-03} before claiming verified offline behaviour.\\
\end{longtable}}

\subsubsection{RS-06 --- Academic Workload and Capacity}
\textbf{Prevent.} Collect availability and assessment dates, split delivery into
small owned tasks, reserve integration/test time and review remaining effort
against actual capacity. Identify backups for dependent work. Do not assume nine
people can all work simultaneously or replace elapsed lead time with person-hours.

\textbf{Respond.} When capacity is insufficient, rebalance tasks and propose a
reviewed scope reduction under \decisionref{Q-SCOPE-01}. Do not silently remove
required permissions, integrity checks or evidence. Record the effect on estimates
and milestones. \textbf{Owner/review:} Project coordinator; weekly and immediately
when a task dependency is at risk.

\subsubsection{RS-01 --- Architecture Validation}
\textbf{Prevent.} Validate the selected stack and API/data contracts before expanding
implementation. Build one minimal path: authenticate, prepare permitted data, save
an identified work proposal durably, reconnect and accept it idempotently in
PostgreSQL. Include a stale-revision case. Record deployment/service choices under
\decisionref{Q-ARCH-01} and \decisionref{Q-ARCH-02}.

\textbf{Respond.} If a service or library blocks that path, evaluate a simpler
compatible implementation with a recorded decision and migration impact. Avoid an
unreviewed stack switch or bypassing central authority. \textbf{Owner/review:}
Backend/architecture lead; next implementation planning session and at each
blocked integration checkpoint until the slice passes.

\subsubsection{RS-07 --- Scope Expansion}
\textbf{Prevent.} Maintain the included/excluded scope in Section~\ref{sec:scope},
maintain the single-final-delivery baseline and estimate additions before
accepting them. Payroll, wage calculations, payments, invoicing and automatic
ticket extraction remain excluded. Keep the agreed cooperative-reported values
and campaign indicators; they are not payroll or accounting.

\textbf{Respond.} Defer optional additions or exchange scope through a reviewed
decision with effort, dependencies and acceptance impact. Do not promise new
CSV/Excel exports as a mitigation; use only outputs already selected for delivery.
\textbf{Owner/review:} Project coordinator; each new request and milestone review.

\subsubsection{RS-08 --- Client Clarification}
\textbf{Prevent.} The client intermediary collects focused questions, requests
sanitised ticket/document examples and schedules concise validation sessions.
Connect each blocking question to Section~\ref{sec:uc-open-decisions}, an owner
and the task that needs the answer. Use synthetic examples while protecting real
worker data.

\textbf{Respond.} If the client is unavailable, record an explicitly provisional
team assumption and develop only reversible parts that do not depend on an
unvalidated business/legal conclusion. Escalate blocking decisions; the
intermediary does not acquire binding product/legal approval authority.
\textbf{Owner/review:} Client intermediary; weekly and before dependent-task dates.

\subsubsection{RS-09 --- Task Coordination}
\textbf{Prevent.} Maintain the shared task board with one accountable owner,
reviewer where needed, due date, dependencies and completion evidence. Record
actual person-hours in the existing shared sheet. Use brief asynchronous updates
for completed work, next work and blockers; create a handover for absences.

\textbf{Respond.} Rebalance overloaded or unowned work in the team review and
update the plan and estimates. Reconcile document contributions before final
report integration. \textbf{Owner/review:} Project coordinator; weekly and whenever
ownership or availability changes.

\subsubsection{RS-02 --- Offline Synchronisation}
\textbf{Prevent.} Implement and test the protocol in Section~\ref{sec:offline-sync}
early: stable record/operation IDs, atomic local saves, ordered dependencies,
idempotent server effects, optimistic revisions and explicit conflict review.
Test normal restart, uncertain outcomes, concurrent starts/edits, worker creation,
authorisation expiry and protected pending data. A transport operation log is not
an append-only restriction on editable business records.

\textbf{Respond.} If automatic replay is unreliable, expose a clear authorised
retry/review action while fixing it. A manual trigger uses the same durable queue,
permissions and conflict validation; it does not introduce Last Write Wins or
discard pending proposals. If necessary, revise final-delivery scope explicitly.
\textbf{Owner/review:} Backend/synchronisation lead; each vertical-slice increment
and before demonstrating verified offline behaviour.

\subsubsection{Responses for the Remaining Risks}
The remaining risks retain monitoring and proportionate actions even when not in
the first six priorities.
{\compacttable
\begin{longtable}{L{0.13}L{0.87}}
\caption{Responses for other monitored risks.}\label{tab:other-risk-responses}\\
\toprule\rowcolor{TableHeader}\textbf{ID} & \textbf{Prevention and response}\\\midrule\endfirsthead
\toprule\rowcolor{TableHeader}\textbf{ID} & \textbf{Prevention and response}\\\midrule\endhead
\multicolumn{2}{r}{\footnotesize\itshape Continued on next page}\\\endfoot
\bottomrule\endlastfoot
\identifier{RS-03} & Review phone/tablet workflows under the approved USA criteria and representative outdoor conditions. If tasks fail, simplify the affected screen and repeat the affected task checks.\\\midrule
\identifier{RS-04} & Establish build and core-rule regression checks before broad integration. Block/review a failing merge, fix the cause and repeat affected regression checks.\\\midrule
\identifier{RS-05} & Require shared contracts, dependency review and peer review for all code regardless of authoring tool. Rework incompatible contributions before merge and record required follow-up.\\\midrule
\identifier{RS-10} & Integrate small changes against reviewed contracts/migrations. Resolve conflicting branches with the relevant owners and rerun affected build/integration checks.\\\midrule
\identifier{RS-11} & Record a backup and handover for changed ownership. Assign an interim owner and recover outstanding decisions/work when the trigger occurs.\\\midrule
\identifier{RS-12} & Obtain sanitised examples and remote walkthroughs on representative devices. If field validation is unavailable, state the limitation and keep dependent assumptions open.\\
\end{longtable}}


\subsection{Monitoring, Ownership and Escalation}
Cec\'{i}lia and Pedro coordinate current risk tracking through Phase 1, drawing
on the assessment contributed by Pepe and Mateo. The role owners in the
register are proposed response responsibilities; assign a named person and backup
before treating an action as scheduled. Review priority risks weekly and the full
register at agreed milestones, scope changes and significant incidents. Escalate
active triggers immediately when they threaten delivery, access or data integrity.

Each review records date, named owner, status (Open, Monitoring, Materialised or
Closed), rating, action/due date, trigger evidence and residual exposure. These are
fields to maintain, not fabricated current results. Materialised risks become issues
with actions; closure needs a reason and evidence. Link mitigation work to the
backlog, remaining estimates and the hours sheet. Reassess exposure after action;
planning mitigation alone does not demonstrate reduced residual risk.

Client questions are coordinated through the intermediary without granting binding
business/legal approval authority. Communicate assumptions and changes at planned
checkpoints, seeking explicit feedback on scope, schedule, resources and quality. Retain detailed votes/action history in the
project artefacts; update the report when the working baseline changes.

\paragraph{Operational examples and requirement risk.}
The examples reviewed on 21 November reduce uncertainty about available forms and
ticket labels. The client-confirmed assigned kg, accepted MEDIANERIA percentage and reported-only penalty policy now resolve the source interpretation. The four supplied Spanish forms are selected; rendering and operational checks remain team tasks. Track these
items against the existing stakeholder/requirements risks and reassess affected
effort before scheduling implementation. Receipt of examples alone does not close
a risk or prove a mitigation completed.

\subsection{Planning and Effort Summary}\label{sec:project-planning}
The supplied hours workbook snapshot of 24 November totals 112 person-hours.
The Summary values reconcile with individual entries plus shared Team meetings
counted for the nine listed members. Two task entries still lack dates;
confirm those dates and any unrecorded work with their authors. The sheet records
past effort and does not establish acceptance of an implementation task.

The supplied estimation workbook contains 56 three-point task estimates with
\(E=(B+4M+W)/6\), totalling 872.05 person-hours. It mixes semester-wide effort and
proof-of-concept assumptions, including older exclusions and 810/950-hour budget
assumptions. The revised 88-task catalogue retains its separate inherited reference
triples, recomputes their PERT values and identifies scope changes/reuse evidence.
The 87 inherited triples total 452.17 hours when full-precision PERT is summed.
The source catalogue total of 453.2 hours summed already-rounded rows. One new
DEV-06 deployment/backup task has a proposed 4/8/16-hour triple (8.67 expected),
bringing reference/proposed catalogue effort to 460.83 hours. All changed scope
and reuse must be re-estimated before these values can become remaining work.
Its task IDs and granularities differ from the 56-item workbook. Neither total is
an approved current remaining-effort forecast. Reconcile task ownership, completed
scope and remaining estimates before calculating actual plus remaining effort;
do not add the 112 hours to either whole-scope historical total or claim a reserve.

\subsubsection{Phase Responsibilities}\label{sec:phase-assignment}
The allocation below follows the supplied team note. Each pair owns coordination
and completion evidence for its phase; choose one execution lead and a reviewer
for each task on the board. Recurring management and early testing continue across
phases. The detailed catalogue supplies task IDs, prerequisites and reference effort.
{\compacttable\begin{longtable}{L{0.12}L{0.5}L{0.38}}
\toprule\rowcolor{TableHeader}\textbf{Phase} & \textbf{Deliverable focus} & \textbf{Responsible pair}\\\midrule\endfirsthead
\toprule\rowcolor{TableHeader}\textbf{Phase} & \textbf{Deliverable focus} & \textbf{Responsible pair}\\\midrule\endhead
\bottomrule\endlastfoot
1 & Requirements, handbook, C4 and API/sync design; management tracking. & Cec\'{i}lia and Pedro.\\\midrule
2 & Existing monorepo/environment/PWA foundation, shared contracts and calculations. & Hugo and Pepe.\\\midrule
3 & Domain Prisma schema/migrations, seed data and local Dexie/outbox structures. & Mateo and Navega.\\\midrule
4 & NestJS API/sync core, protected files, account provisioning/authentication PWA infrastructure, private deployment and backups. & Andr\'{e} and Gabriel.\\\midrule
5 & Agricultural backend/UI modules, manager delegation, optional planning, documents and indicators. & Pepe and Dinis.\\\midrule
6 & Integrated QA, client walkthrough, deployment/handover and final academic deliverables. & Gabriel and Navega.\\
\end{longtable}}
Phase numbers group responsibility and dependencies; readiness is assessed per
task/contract. Integration waits for its listed prerequisites and validation;
independent preparation can proceed in parallel. One final delivery is planned,
with final calendar and capacity unconfirmed under \decisionref{Q-PM-01}. Pepe,
Gabriel and Navega participate in two phases, so schedule their assignments
against actual availability. Phase pairs are confirmed by the supplied note;
individual task, risk, test and operational leads, dates and team-wide AI inventory
still need assignment under \decisionref{Q-PM-02}.

